%%%%%%%%%%%%%%%%%%%%%%%%%%%%%%%%%%%%%%%%%%%%%%%%%%%%%%%%%%%%%%%%%%%%%%%%%%%
% Función: capítulo de ejemplo de conclusiones de la tesis por compendio.
%
% Uso: modifica este fichero para adaptar el contenido a tu trabajo.
%%%%%%%%%%%%%%%%%%%%%%%%%%%%%%%%%%%%%%%%%%%%%%%%%%%%%%%%%%%%%%%%%%%%%%%%%%%

\ifthenelse{\equal{\myLanguage}{english}}
{
  \chapter{Conclusions and Future Work}
}
{
  \chapter{Conclusiones y líneas futuras}
}
\label{cha:compendium-conclusions}

\ifthenelse{\equal{\myLanguage}{english}}
{
  For an example of connecting conclusions to scope, compare the application-level goals in functional assessment~\cite{nieva2024proposal}, the system context described for EYEFUL~\cite{nieva2023eyeful} and the different sensing setting of vessel localisation~\cite{ramirez2026jstars}. Use the comparison to ask which claims are supported by your own evidence and which would require further validation. Do not present results from these publications as results of this fictional thesis.
}
{
  Como ejemplo de conexión entre conclusiones y alcance, compara los objetivos ligados a la aplicación en evaluación funcional~\cite{nieva2024proposal}, el contexto del sistema descrito para EYEFUL~\cite{nieva2023eyeful} y el escenario de sensado diferente de la localización de barcos~\cite{ramirez2026jstars}. Utiliza la comparación para preguntarte qué afirmaciones están sustentadas por tus propias evidencias y cuáles exigirían más validación. No presentes resultados de esas publicaciones como resultados de esta tesis ficticia.
}

\ifthenelse{\equal{\myLanguage}{english}}
{
  \section{Main Conclusions}
}
{
  \section{Conclusiones principales}
}

\ifthenelse{\equal{\myLanguage}{english}}
{
  Summarize the conclusions supported by the thesis as a whole, relate them to the objectives and discuss the principal limitations. Close the extended summary by identifying research directions that follow from the results.
}
{
  Resume las conclusiones sustentadas por el conjunto de la tesis, relaciónalas con los objetivos y discute las principales limitaciones. Cierra el resumen extendido identificando las líneas de investigación que se derivan de los resultados.
}

\ifthenelse{\equal{\myLanguage}{english}}
{
  \section{Limitations}

  Discuss the assumptions and practical constraints that limit the generalization of the results and should be considered when interpreting the conclusions.

  \section{Future Research Directions}
  \subsection{Immediate extensions}

  Identify improvements and experiments that follow directly from the work presented in the thesis.

  \subsection{Long-term research opportunities}

  Outline broader questions opened by the combined contributions and publications.
}
{
  \section{Limitaciones}

  Discute las hipótesis y restricciones prácticas que limitan la generalización de los resultados y deben considerarse al interpretar las conclusiones.

  \section{Líneas futuras de investigación}
  \subsection{Extensiones inmediatas}

  Identifica las mejoras y experimentos que se derivan directamente del trabajo presentado en la tesis.

  \subsection{Oportunidades de investigación a largo plazo}

  Plantea cuestiones más amplias abiertas por el conjunto de las contribuciones y publicaciones.
}

%%% Local Variables:
%%% TeX-master: "../../book"
%%% End:
